\section{An obstruction to monotonicity on either parity}
\label{sec:parity}

At a fixed nonreal zero the leading quadrature error has a rotating
complex phase. Its projection onto the imaginary part of the root
displacement determines the signed change of the counted discriminant
\(D\) \cite[\S 2.2, Eq.~(2.4)]{agc}. The compact-source construction
below uses this phase to test monotonicity in the truncation
degree \cite{jh05}.

\Needspace{15\baselineskip}
\begin{proposition}[Displacement for a compact source]
\label{p:displacement}
Let \(\mu\) be an even probability measure with infinite support in
\([-K,K]\), and take \(F(z)=\int e^{izu}\,d\mu(u)\) and its
\(N\)-site Jacobi readout \(G_N\). Suppose a fixed circle centered on
the real axis encloses exactly two zeros of \(F\), a simple pair
\(\zeta=x+iy\) and \(\bar\zeta\) with \(x,y>0\), and that its boundary contains no zero.
For all sufficiently large \(N\), it encloses a simple pair
\(\zeta_N,\bar\zeta_N\) of \(G_N\), with \(\zeta_N\to\zeta\)
and a common zero-free collar about the circle. Put
\[
 A_N=\frac{h_N^*|\zeta|^{2N}}{(2N)!},\qquad
 \theta=\arg\zeta,\qquad D_N=-(\operatorname{Im}\zeta_N)^2 .
\]
Then the displacement and the two-degree increment are
\begin{align}
 \zeta_N-\zeta
   &=\frac{(-1)^N h_N^*\zeta^{2N}}{(2N)!\,F'(\zeta)}
     +O(A_N/N), \label{p:root-shift}\\
 \frac{D_{N+2}-D_N}{A_N}
   &=2y(-1)^N\operatorname{Im}
      \left(\frac{e^{2iN\theta}}{F'(\zeta)}\right)+O(N^{-1}).
      \label{p:degree-sign}
\end{align}
The implied constants are uniform in \(N\) with the source and
circle fixed.
\end{proposition}

\begin{proof}
The moment Gram matrices are positive definite by infinite support.
Multiplication by \(u\) on \(L^2(\mu)\) has norm at most \(K\), so
its finite Jacobi compressions satisfy \(\|J_N\|\le K\) and the
off-diagonal links satisfy \(0<\beta_j\le K^2\). In particular
\(0<h_N^*=\prod_{j=1}^N\beta_j\le K^{2N}\).

All interpolation knots lie in \([-K,K]\), so the Gaussian quadrature
error representation \eqref{eq:full-phase} applies on every fixed
complex disk. Write
\[
 H_N(z):=F(z)-G_N(z)=L_N(z)E_N(z),\qquad
 L_N(z)=(-1)^N\frac{h_N^*z^{2N}}{(2N)!},
 \qquad E_N(z)=\mathbb E e^{izX_N}.
\]
Here \(X_N\) is a uniform simplex average of the doubled Gauss
nodes and an independent source point \(U\) of law
\(P_N(u)^2\,d\mu(u)/h_N^*\), where \(P_N\) is monic. Reflection
makes \(X_N\) symmetric and \(|X_N|\le K\).
For the \(m=2N+1\) simplex coordinates \(S_j\), the identity
\(\mathbb E(S_j S_k)=(1+\mathbf 1_{j=k})/[m(m+1)]\) gives
\[
 \mathbb E X_N^2
   =\frac{2(\mathbb E U^2+\operatorname{tr} J_N^2)}{(2N+1)(2N+2)}
   \le \frac{K^2}{2N+1},
\]
since the doubled nodes sum to zero and have total squared sum
\(2\operatorname{tr}J_N^2\) \cite[\S 7, second moment]{vcd}.
Taylor's integral remainder, with
\(\mathbb E X_N=0\), now yields
\begin{equation}
 |E_N(z)-1|
 \le \frac{K^2|z|^2 e^{K|\operatorname{Im} z|}}{2(2N+1)},\qquad
 \sup_{|z|\le R}|H_N(z)|
 \le U_N(R):=\frac{e^{KR}h_N^*R^{2N}}{(2N)!}.
 \label{p:compact-bound}
\end{equation}

Choose disjoint closed isolating disks around the two zeros inside
the circle, arranged in conjugate positions. Choose \(R\) to
cover the circle, a closed zero-free collar of \(F\) and a fixed
positive margin around the disks. Since \(U_N(R)\to0\), Rouch\'e's
theorem gives counts \(2,1,1\) on these three circles, and uniform
convergence gives the common cleared collar. Let \(d_N=\zeta_N-\zeta\).
The positive minimum of \(|F(z)/(z-\zeta)|\) on the upper disk
gives \(d_N=O(U_N(R))\), using
\(F(\zeta_N)=H_N(\zeta_N)\). The fixed margin and Cauchy's
estimate bound \(H_N'\) on the disk by a constant times \(U_N(R)\).
Taylor expansion of \(F\) and integration of \(H_N'\) along the
displacement therefore give
\[
 F'(\zeta)d_N=L_N(\zeta)E_N(\zeta)+O(U_N(R)^2).
\]
The absolute disk bound retains the squared norm \(h_N^*\).
After division by the root scale it gives
\[
 \frac{U_N(R)^2}{A_N}
   =\frac{e^{2KR}h_N^*}{(2N)!}
       \left(\frac{R^2}{|\zeta|}\right)^{2N}
   \le \frac{e^{2KR}}{(2N)!}
       \left(\frac{KR^2}{|\zeta|}\right)^{2N}
   =o(N^{-1}).
\]
Together with \eqref{p:compact-bound} this proves
\eqref{p:root-shift}. Expanding \(D_N=-(y+\operatorname{Im}d_N)^2\)
and using \(A_N=o(N^{-1})\) gives
\[
 D_N+y^2=-2y(-1)^N A_N\operatorname{Im}
   \left(\frac{e^{2iN\theta}}{F'(\zeta)}\right)+O(A_N/N).
\]
Furthermore,
\[
 \frac{A_{N+2}}{A_N}
  =\frac{\beta_{N+1}\beta_{N+2}|\zeta|^4}
   {(2N+1)(2N+2)(2N+3)(2N+4)}
  \le \frac{K^4|\zeta|^4}{(2N+1)(2N+2)(2N+3)(2N+4)} .
\]
Thus \(D_{N+2}+y^2=O(A_NN^{-4})\), which proves
\eqref{p:degree-sign} and fixes the sign of its leading term.
\end{proof}

\begin{theorem}[Both signs occur on both parities]
\label{p:oscillation}
Set \(T(u)=(1-|u|)_+\) and take the continuous even density
\begin{equation}
 d\mu(u)=\frac{T(u)+[T(u-1)+T(u+1)]/4}{3/2}\,du .
 \label{p:density}
\end{equation}
Let \(D_N\) be the discriminant of the Jacobi pair counted inside
\(\lvert z-\pi\rvert=3/2\) for all sufficiently large \(N\). Put
\(a=\operatorname{arcosh}2\) and
\(A_N=h_N^*(\pi^2+a^2)^N/(2N)!\). For each \(p\in\{0,1\}\),
\begin{align}
 \limsup_{\substack{N\to\infty\\N\equiv p\pmod 2}}
     \frac{D_{N+2}-D_N}{A_N}
   &\ge \frac{a(\pi^2+a^2)}{2\sqrt{3}}, \notag\\
 \liminf_{\substack{N\to\infty\\N\equiv p\pmod 2}}
     \frac{D_{N+2}-D_N}{A_N}
   &\le -\frac{a(\pi^2+a^2)}{2\sqrt{3}} .
 \label{p:limbounds}
\end{align}
Consequently each parity contains infinitely many strict increases
and infinitely many strict decreases of \(D_N\).
\end{theorem}

\begin{proof}
Each tent has integral one. The density is positive on \((-2,2)\)
and has support \([-2,2]\). Integrating the central tent and
translating the two others gives
\[
 F(z)=\left[\frac{\sin(z/2)}{z/2}\right]^2
       \frac{1+\cos(z)/2}{3/2}.
\]
The squared sinc has zeros at the nonzero multiples of \(2\pi\).
The other factor has simple zeros at \((2k+1)\pi\pm ia\),
\(k\in\mathbb Z\). Since
\(\cosh(3/2)>1+(3/2)^2/2>2\), the two zeros
\(\zeta=\pi+ia\) and \(\bar\zeta\) lie strictly inside the
circle. Each other zero has horizontal distance at least \(\pi\)
from its center. In particular Proposition~\ref{p:displacement}
applies and \(0<4\theta=4\arctan(a/\pi)<\pi\).

At \(\zeta\), \(\sin\zeta=-i\sqrt{3}\) and
\(\sin^2(\zeta/2)=\cosh^2(a/2)=3/2\), giving
\(F'(\zeta)=2i\sqrt{3}/\zeta^2\). Formula \eqref{p:degree-sign}
becomes
\[
 \frac{D_{N+2}-D_N}{A_N}
   =-(-1)^N\frac{a(\pi^2+a^2)}{\sqrt{3}}
      \cos\bigl(2(N+1)\theta\bigr)+O(N^{-1}).
\]
On \(N=p+2j\) the normalized leading term divided by
\(a(\pi^2+a^2)/\sqrt{3}\) is a sinusoid \(b_j\) with angular step
\(4\theta\). Its first two Ces\`aro moments are \(0\) and \(1/2\).
Indeed the finite geometric sums at frequencies \(4\theta\) and
\(8\theta\) are bounded because neither frequency is a multiple of
\(2\pi\). For \(B=\limsup b_j\) and \(\epsilon>0\), eventually
\((b_j+1)(B+\epsilon-b_j)\ge0\), so
\[
 b_j^2\le B+\epsilon+(B+\epsilon-1)b_j .
\]
Taking Ces\`aro means and letting \(\epsilon\) decrease to zero
gives \(B\ge1/2\). The same argument for \(-b_j\) gives
\(\liminf b_j\le-1/2\). This proves \eqref{p:limbounds}
for both parities, with \(A_N>0\).
\end{proof}

The example excludes an eventual monotonicity inference based on
positivity, evenness, the cosine observation and a fixed count-two
chart for continuous compactly supported sources at time zero. The
canonical problem retains the full theta profile and its square
arithmetic \cite[\S 1.2, Eq.~(1.1)]{agc}.
